% Insert inside the shared Literature Review section; Member 3 adds core/datapath studies.
% Merge member2_references.bib into the team's bibliography without duplicating keys.
\subsection{RISC Principles and the RISC-V ISA}
Reduced instruction set computer designs use regular instructions and a
load/store organization to make processor implementations easier to reason
about. RV32I follows this organization: arithmetic operates on registers,
while load and store instructions transfer data between registers and memory.
Its base instructions are 32 bits long, and source and destination register
fields occupy consistent positions across formats to simplify decoding.
RV32I provides 32 registers of 32 bits each, with register x0 fixed at
zero~\cite{riscv_unpriv_2026}.

Asanovi\'{c} and Patterson argue that an instruction set is an interface
specification, distinct from a particular processor implementation. Their
case for an open ISA emphasizes research, education, and a small base
instruction set with optional extensions~\cite{asanovic_patterson_2014}.
This distinction permits different microarchitectures to implement the
same instruction semantics; selecting RISC-V does not itself require a
single-cycle datapath.

The current RV32I specification defines ADD and SUB as operations whose
low 32 result bits are retained, with arithmetic overflow ignored. AND, OR,
and XOR operate bitwise. BEQ and BNE compare registers for equality and
inequality, respectively~\cite{riscv_unpriv_2026}. These semantics motivate
the project's 32-bit ALU and its subtraction-based equality check. The
ALU's zero signal is an internal datapath connection, rather than an
architectural condition-code register.

The educational need addressed here is a small, inspectable implementation
that connects instruction semantics to modular RTL and repeatable tests.
The project targets ten instructions: ADD, SUB, AND, OR, XOR, ADDI, LW,
SW, BEQ, and BNE. This is a teaching subset of RV32I; it does not claim
complete RV32I conformance. Member 3's review of existing datapaths and
cores should complete the comparison with prior implementations.
